\documentclass[12pt,a4paper]{article}
\usepackage[utf8]{inputenc}
\usepackage[T1]{fontenc}
\usepackage[brazilian]{babel}
\usepackage{amsmath,amssymb}
\usepackage[margin=2.2cm]{geometry}
\usepackage{tikz}
\usetikzlibrary{arrows.meta,positioning,decorations.pathreplacing}
\usepackage{tcolorbox}
\tcbuselibrary{skins,breakable}
\usepackage{enumitem}
\usepackage{xcolor}
\usepackage{booktabs}
\usepackage{hyperref}
\hypersetup{colorlinks=true,linkcolor=blue!50!black,urlcolor=blue!50!black}

\newtcolorbox{definicao}[1]{colback=blue!5,colframe=blue!55!black,fonttitle=\bfseries,coltitle=white,title={#1},breakable,boxrule=0.8pt,arc=2mm}
\newtcolorbox{exemplo}[1]{colback=green!5,colframe=green!35!black,fonttitle=\bfseries,coltitle=white,title={#1},breakable,boxrule=0.8pt,arc=2mm}
\newtcolorbox{atencao}[1]{colback=red!5,colframe=red!55!black,fonttitle=\bfseries,coltitle=white,title={#1},breakable,boxrule=0.8pt,arc=2mm}
\newtcolorbox{dica}[1]{colback=yellow!12,colframe=orange!70!black,fonttitle=\bfseries,coltitle=black,title={#1},breakable,boxrule=0.8pt,arc=2mm}
\newtcolorbox{macete}[1]{colback=purple!5,colframe=purple!55!black,fonttitle=\bfseries,coltitle=white,title={#1},breakable,boxrule=0.8pt,arc=2mm}

\title{\vspace{-1.5cm}\textbf{Tema 0 — Fundamentos Aritméticos e Macetes}\\[2pt]\large A base que sustenta tudo: sinais, ordem das operações e frações}
\author{}
\date{}

\begin{document}
\maketitle
\thispagestyle{empty}

\section*{Por que este tema vem antes de tudo}
Todo erro em equação, inequação, potência ou função quase sempre não é
``erro de matéria nova'' — é erro de \textbf{conta básica}: um sinal
trocado, uma ordem de operação ignorada, uma fração mal somada. Domine
este arquivo e a taxa de erro boba cai demais.

\section{Regra de sinais — a tabela mais importante da sua vida matemática}

\subsection{Adição e subtração de números com sinal}

\begin{dica}{A ideia da reta numérica: andar para a direita ou esquerda}
Pense em cada número como um passo numa reta. Número \textbf{positivo}:
ande para a \textbf{direita}. Número \textbf{negativo}: ande para a
\textbf{esquerda}. Somar é continuar andando na mesma direção do sinal;
\textbf{subtrair um número é andar na direção oposta à dele}.
\end{dica}

\begin{center}
\begin{tikzpicture}[scale=1]
  \draw[-{Stealth[length=2.5mm]}, thick] (-5.5,0) -- (5.5,0);
  \foreach \x in {-5,-4,-3,-2,-1,0,1,2,3,4,5}
    \draw (\x,0.08) -- (\x,-0.08) node[below]{\small \x};
  \draw[thick, blue, -{Stealth[length=3mm]}] (0,0.6) to[bend left=15] (3,0.6);
  \node[above,blue] at (1.5,0.9) {\small $+3$ (direita)};
  \draw[thick, red, -{Stealth[length=3mm]}] (3,-0.6) to[bend left=15] (1,-0.6);
  \node[below,red] at (2,-0.9) {\small $-2$ (esquerda)};
  \filldraw[black] (0,0) circle(2pt);
  \filldraw[black] (1,0) circle(2pt);
\end{tikzpicture}
\end{center}
No desenho: começando em $0$, ando $3$ para direita (chego no $3$), depois
ando $2$ para a esquerda (chego no $1$). Isso é $0+3-2=1$.

\begin{definicao}{Regra prática para somar/subtrair dois números com sinal}
\begin{itemize}[nosep]
  \item \textbf{Sinais iguais}: some os valores (ignorando o sinal) e mantenha o sinal comum.
  \[
  (-4)+(-7) = -(4+7) = -11
  \]
  \item \textbf{Sinais diferentes}: subtraia o menor valor do maior (ignorando sinais) e
  fique com o sinal do número que tem \textbf{maior valor absoluto}.
  \[
  (-9)+5 = -(9-5) = -4 \qquad\qquad 9+(-5) = +(9-5)=4
  \]
\end{itemize}
\end{definicao}

\begin{atencao}{O macete que resolve 90\% dos erros de subtração: ``menos com menos vira mais''}
$a - (-b) = a+b$. Subtrair um número negativo é o mesmo que somar o
positivo correspondente — os dois sinais de menos ``se cancelam''.
\[
7 - (-3) = 7+3 = 10 \qquad\qquad -2-(-5) = -2+5 = 3
\]
Sempre que vir \texttt{menos, abre parêntese, menos}, transforme os dois
sinais em um \texttt{+} \textbf{antes} de fazer qualquer conta.
\end{atencao}

\subsection{Multiplicação e divisão de números com sinal}

\begin{macete}{O ``jogo da velha'' dos sinais — decore essa tabela de olhos fechados}
\begin{center}
\renewcommand{\arraystretch}{1.6}
\begin{tabular}{c|cc}
$\times$ ou $\div$ & \textbf{mesmo sinal} & \textbf{sinais diferentes}\\
\hline
\textbf{resultado} & \textcolor{blue}{$+$} & \textcolor{red}{$-$}
\end{tabular}
\end{center}
\[
(+)\cdot(+) = + \qquad (-)\cdot(-) = + \qquad (+)\cdot(-) = - \qquad (-)\cdot(+) = -
\]
Frase para lembrar: \emph{``amigo do meu amigo é meu amigo (+); inimigo do
meu inimigo é meu amigo (+); amigo do meu inimigo é meu inimigo (-)''}.
Vale exatamente igual para \textbf{divisão}.
\end{macete}

\begin{exemplo}{Aplicando}
\[
(-3)\cdot(-4) = 12 \qquad\qquad (-3)\cdot 4 = -12 \qquad\qquad \frac{-20}{-5}=4 \qquad\qquad \frac{-20}{5}=-4
\]
\end{exemplo}

\begin{atencao}{Cuidado: regra de sinal de multiplicação NÃO vale para soma}
$(-3)+(-4) = -7$ (soma sinais iguais $\to$ mantém o sinal e \textbf{soma} os
valores), mas $(-3)\cdot(-4)=+12$ (multiplicação sinais iguais $\to$
\textbf{positivo}). São regras diferentes para operações diferentes —
muita gente mistura as duas.
\end{atencao}

\section{Potência com sinal — e por que ``subtração de números elevados'' trava tanta gente}

\begin{atencao}{Regra \#1: o parêntese decide se o sinal entra na potência}
\[
(-2)^2 = (-2)\cdot(-2) = 4 \qquad\qquad\qquad -2^2 = -(2^2) = -4
\]
Sem parêntese, o expoente age só no número, e o sinal de menos fica
``do lado de fora'', sendo aplicado só no final.
\end{atencao}

\begin{definicao}{Expoente par ou ímpar decide o sinal do resultado (quando há parêntese)}
\begin{itemize}[nosep]
  \item Expoente \textbf{par}: resultado sempre \textbf{positivo}, não importa o sinal da base.
  $(-3)^2=9,\ (-3)^4=81$.
  \item Expoente \textbf{ímpar}: resultado \textbf{mantém} o sinal da base.
  $(-3)^3=-27,\ (-3)^5=-243$.
\end{itemize}
\end{definicao}

\begin{atencao}{O erro mais caro da prova: $a^2-b^2 \neq (a-b)^2$}
Isso \textbf{não} é sobre sinal, é sobre a operação em si — são contas
diferentes. Compare com números:
\[
5^2-3^2 = 25-9=16 \qquad\qquad\qquad (5-3)^2 = 2^2=4
\]
$16 \neq 4$! O que \emph{é} verdade (Tema 4 — produtos notáveis) é
$a^2-b^2=(a-b)(a+b)$, que é diferente de $(a-b)^2$. Não confunda
\textbf{subtrair dois quadrados} com \textbf{elevar uma subtração ao
quadrado}.
\end{atencao}

\begin{exemplo}{``Subtraindo números elevados'' — dois cenários bem diferentes}
\textbf{Cenário A — só números, calcule cada potência primeiro:}
\[
5^3 - 2^3 = 125-8=117
\]
\textbf{Cenário B — termos com $x$: só pode ``juntar'' potências de mesma
base e mesmo expoente (termos semelhantes):}
\[
5x^3-2x^3 = 3x^3 \qquad\qquad\text{mas}\qquad\qquad 5x^3-2x^2 \ \text{já não simplifica mais}
\]
\end{exemplo}

\begin{macete}{Checklist rápido antes de subtrair algo com expoente}
\begin{enumerate}[nosep]
  \item Calcule cada potência \textbf{separadamente} primeiro (respeitando parênteses).
  \item Só depois disso, aplique a regra de sinais da soma/subtração.
  \item Se são termos algébricos ($x^n$), só combine quando a \textbf{base e o expoente} forem idênticos.
\end{enumerate}
\end{macete}

\section{Ordem das operações — a hierarquia que ninguém pode ignorar}

\begin{definicao}{A hierarquia (de cima para baixo, sempre nessa ordem)}
\begin{center}
\begin{tikzpicture}[
  every node/.style={draw,rounded corners,minimum width=8.5cm,minimum height=0.9cm,font=\small},
  node distance=2mm
]
  \node[fill=red!15] (n1) {\textbf{1º} — Parênteses, colchetes e chaves: $(\ )\ [\ ]\ \{\ \}$};
  \node[fill=orange!15, below=of n1] (n2) {\textbf{2º} — Potências e raízes: $x^n,\ \sqrt[n]{x}$};
  \node[fill=yellow!25, below=of n2] (n3) {\textbf{3º} — Multiplicação e divisão (esquerda $\to$ direita)};
  \node[fill=green!15, below=of n3] (n4) {\textbf{4º} — Adição e subtração (esquerda $\to$ direita)};
  \draw[-{Stealth[length=3mm]},thick] (n1) -- (n2);
  \draw[-{Stealth[length=3mm]},thick] (n2) -- (n3);
  \draw[-{Stealth[length=3mm]},thick] (n3) -- (n4);
\end{tikzpicture}
\end{center}
\end{definicao}

\begin{atencao}{Multiplicação NÃO vem antes de divisão — elas empatam}
Multiplicação e divisão têm a \textbf{mesma prioridade} entre si; resolve-se
na ordem em que aparecem, da esquerda para a direita. O mesmo vale para
soma e subtração entre si. Não existe ``multiplicação sempre antes de
divisão''.
\end{atencao}

\begin{exemplo}{Um exemplo que usa a hierarquia inteira, passo a passo}
\[
20 - 2\cdot(3+1)^2 \div 4
\]
\textbf{Passo 1 (parênteses)}: $3+1=4$.
\[
20 - 2\cdot 4^2 \div 4
\]
\textbf{Passo 2 (potência)}: $4^2=16$.
\[
20 - 2\cdot 16 \div 4
\]
\textbf{Passo 3 (multiplicação/divisão, esquerda para direita)}: $2\cdot16=32$, depois $32\div4=8$.
\[
20 - 8
\]
\textbf{Passo 4 (subtração)}:
\[
= 12
\]
\end{exemplo}

\begin{macete}{Truque para nunca se perder em conta grande}
Reescreva a expressão a cada passo, mudando \textbf{só} o pedaço que você
acabou de calcular — como fizemos acima. Nunca tente resolver ``de
cabeça'' pulando 2 passos ao mesmo tempo; é aí que mora o erro bobo.
\end{macete}

\section{Frações — as 4 operações}

\begin{definicao}{Soma e subtração}
\textbf{Mesmo denominador}: opera só os numeradores e repete o denominador.
\[
\frac{3}{7}+\frac{2}{7} = \frac{5}{7}
\]
\textbf{Denominadores diferentes}: primeiro reduza ao mesmo denominador
(usando o MMC, seção abaixo), depois some/subtraia os numeradores.
\[
\frac{1}{4}+\frac{1}{6} = \frac{3}{12}+\frac{2}{12} = \frac{5}{12}
\]
\end{definicao}

\begin{center}
\begin{tikzpicture}[scale=0.85]
  % pizza 1: 1/4
  \begin{scope}
    \draw[thick] (0,0) circle (1.2cm);
    \draw[thick] (0,0) -- (90:1.2);
    \draw[thick] (0,0) -- (0:1.2);
    \fill[blue!30] (0,0) -- (0:1.2) arc (0:90:1.2) -- cycle;
    \node at (45:0.6) {\small $\frac14$};
    \node[below] at (0,-1.4) {$\frac14 = \frac{3}{12}$};
  \end{scope}
  % pizza 2: 1/6
  \begin{scope}[xshift=3.2cm]
    \draw[thick] (0,0) circle (1.2cm);
    \draw[thick] (0,0) -- (90:1.2);
    \draw[thick] (0,0) -- (30:1.2);
    \fill[green!30] (0,0) -- (30:1.2) arc (30:90:1.2) -- cycle;
    \node at (60:0.7) {\small $\frac16$};
    \node[below] at (0,-1.4) {$\frac16 = \frac{2}{12}$};
  \end{scope}
  % pizza 3: soma = 5/12
  \begin{scope}[xshift=6.4cm]
    \draw[thick] (0,0) circle (1.2cm);
    \draw[thick] (0,0) -- (90:1.2);
    \draw[thick] (0,0) -- (-60:1.2);
    \fill[orange!30] (0,0) -- (-60:1.2) arc (-60:90:1.2) -- cycle;
    \node[below] at (0,-1.4) {$\frac{5}{12}$};
  \end{scope}
\end{tikzpicture}
\end{center}

\begin{definicao}{Multiplicação — a mais simples}
Multiplique numerador por numerador, denominador por denominador. Não
precisa de denominador comum!
\[
\frac{2}{3}\cdot\frac{5}{7} = \frac{2\cdot5}{3\cdot7} = \frac{10}{21}
\]
\end{definicao}

\begin{definicao}{Divisão — multiplique pela fração invertida}
\[
\frac{a}{b}\div\frac{c}{d} = \frac{a}{b}\cdot\frac{d}{c}
\]
\end{definicao}

\begin{exemplo}{$\dfrac{2}{3}\div\dfrac{4}{9}$}
\[
\frac{2}{3}\cdot\frac{9}{4} = \frac{18}{12} = \frac{3}{2}
\]
(simplificamos dividindo $18$ e $12$ por $6$ — veja MDC abaixo).
\end{exemplo}

\begin{macete}{``Copia, inverte e multiplica'' — e por que funciona}
Dividir por uma fração é perguntar ``quantas vezes ela cabe?''. Dividir por
$\frac{c}{d}$ é o mesmo que multiplicar por $\frac{d}{c}$ porque
$\frac{c}{d}\cdot\frac{d}{c}=1$ — você está multiplicando por um jeito
disfarçado de $1$ para ``desfazer'' a fração de baixo.
\end{macete}

\section{MMC e MDC — a ferramenta por trás do denominador comum}

\begin{definicao}{MMC (Mínimo Múltiplo Comum) — para somar/subtrair frações}
Fatore cada número em primos e pegue \textbf{cada primo com o maior expoente} que aparece.
\end{definicao}

\begin{exemplo}{MMC de $12$ e $18$}
\[
12 = 2^2\cdot3 \qquad\qquad 18=2\cdot3^2
\]
Maior expoente de $2$: $2^2$. Maior expoente de $3$: $3^2$.
\[
\text{MMC} = 2^2\cdot3^2 = 4\cdot9=36
\]
\end{exemplo}

\begin{definicao}{MDC (Máximo Divisor Comum) — para simplificar frações}
Fatore em primos e pegue \textbf{cada primo comum com o menor expoente}.
\end{definicao}

\begin{exemplo}{MDC de $18$ e $12$ (para simplificar $\frac{18}{12}$)}
\[
18=2\cdot3^2 \qquad 12=2^2\cdot3 \qquad\Rightarrow\qquad \text{MDC}=2\cdot3=6
\]
\[
\frac{18}{12} = \frac{18\div6}{12\div6} = \frac{3}{2}
\]
\end{exemplo}

\section*{Checklist final — os erros que mais custam pontos}
\begin{itemize}[nosep]
  \item $-3-(-5)$: transforme em $-3+5$ \textbf{antes} de calcular. \emph{(Resposta: $2$.)}
  \item $(-2)^2 = 4$, mas $-2^2=-4$. Confira sempre se tem parêntese.
  \item Sinais iguais na \textbf{soma}: soma valores e mantém sinal. Sinais iguais na \textbf{multiplicação}: dá sempre positivo. São regras diferentes!
  \item $a^2-b^2 \neq (a-b)^2$ — calcule as potências antes de subtrair, a menos que esteja fatorando.
  \item Ordem das operações: parênteses $\to$ potência/raiz $\to$ multiplicação/divisão (esquerda-direita) $\to$ soma/subtração (esquerda-direita).
  \item Fração: só soma/subtrai numerador direto se o denominador já for igual — senão, MMC primeiro.
\end{itemize}

\end{document}
